\section*{Capítulo \ref{ch:b3:modules} --- Módulos sobre um domínio de ideais
principais}

\begin{solution}{exo:b3:modules:1}
(a) Se $\Q = \Z q_1 + \dots + \Z q_k$, seja $d$ um denominador comum dos $q_i$: toda combinação
está em $\frac1d\Z$, mas $\frac1{2d} \notin \frac1d\Z$. Contradição.

(b) Sem \omterm{def:b3:modules:torsion}{torção}: $nq = 0$ com $n \neq 0$ força $q = 0$ em $\Q$. Não
\omterm{def:b3:modules:free}{livre}: dois racionais não nulos $\frac ab, \frac cd$ quaisquer satisfazem a relação não
trivial $(bc)\frac ab - (ad)\frac cd =
0$, de modo que uma base tem no máximo um elemento; $\Q \cong \Z$
tornaria $\Q = \Z q$ cíclico, mas $\frac q2 \notin \Z q$. (E $\Q \neq
0$.)

(c) O~\cref{thm:b3:modules:structure} supõe a \emph{geração finita}, que (a) nega: não há
contradição --- ao contrário, $\Q$ mostra que a hipótese é necessária
no enunciado ``sem \omterm{def:b3:modules:torsion}{torção} $\Rightarrow$ \omterm{def:b3:modules:free}{livre}''.
\end{solution}

\begin{solution}{exo:b3:modules:2}
$360 = 2^3\cdot3^2\cdot5$. Partições: de $3$: $(3), (2,1),
(1,1,1)$; de $2$: $(2), (1,1)$; de $1$: $(1)$.
Logo, $3 \times 2
\times 1 = 6$ grupos. \omterm{thm:b3:modules:structure}{Divisores elementares} $\to$ \omterm{thm:b3:modules:smith}{fatores invariantes}:
\begin{center}
\begin{tabular}{l l}
$\Z/8 \times \Z/9 \times \Z/5$ & $\Z/360$\\
$\Z/8 \times \Z/3 \times \Z/3 \times \Z/5$ & $\Z/3 \times
\Z/120$\\
$\Z/4 \times \Z/2 \times \Z/9 \times \Z/5$ & $\Z/2 \times
\Z/180$\\
$\Z/4 \times \Z/2 \times \Z/3 \times \Z/3 \times \Z/5$ & $\Z/6
\times \Z/60$\\
$(\Z/2)^3 \times \Z/9 \times \Z/5$ & $\Z/2 \times \Z/2 \times
\Z/90$\\
$(\Z/2)^3 \times \Z/3 \times \Z/3 \times \Z/5$ & $\Z/2 \times
\Z/6 \times \Z/30$
\end{tabular}
\end{center}
(Para passar aos \omterm{thm:b3:modules:smith}{fatores invariantes}: o maior $d_s$ reúne a maior
potência de cada primo, e assim por diante, para baixo.) De ordem
$p^5$: tantos quantas as partições de $5$, a saber $7$.
\end{solution}

\begin{solution}{exo:b3:modules:3}
$B$: $D_1 = \gcd(2,4,6,8) = 2$; $D_2 = \abs{\det B} = \abs{16 -
24} = 8$. \omterm{thm:b3:modules:smith}{Fatores invariantes} $d_1 = 2$, $d_2 = 8/2 = 4$: forma de Smith $\operatorname{diag}(2, 4)$,
e $\Z^2/B\Z^2 \cong \Z/2\Z
\times \Z/4\Z$.

$C$: diagonal, mas \emph{não} de Smith ($2 \nmid 3$). $D_1 =
\gcd(2,3,12) = 1$; $D_2 = \gcd(2\cdot3,\, 2\cdot12,\, 3\cdot12) =
\gcd(6, 24, 36) = 6$; $D_3 = 72$. Logo
$d = (1, 6, 12)$ e $\Z^3/C\Z^3 \cong \Z/6\Z\times\Z/12\Z$ --- coerentemente com o teorema chinês dos restos: $\Z/2\times\Z/3\times\Z/12 \cong \Z/6\times\Z/12$.
\end{solution}

\begin{solution}{exo:b3:modules:4}
Escreva $B = Q\,\operatorname{diag}(d_1, \dots, d_n)\,P$ com $Q, P
\in GL_n(\Z)$ (\cref{thm:b3:modules:smith}; nenhum $d_i$ é nulo, pois
$\det B \neq 0$). Então $\Z^n/B\Z^n \cong \Z^n/
\operatorname{diag}(d)\Z^n = \prod_i \Z/d_i\Z$ (o isomorfismo composto $x \mapsto Q^{-1}x$ de $\Z^n$ leva
$B\Z^n$ sobre $\operatorname{diag}(d)P\Z^n = \operatorname{diag}(d)\Z^n$). Sua cardinalidade é $\prod \abs{d_i} = \abs{\det
\operatorname{diag}(d)} = \abs{\det B}$, pois $\det Q, \det P = \pm
1$. Para $L = \Z(2,0) + \Z(1,3)$:
$B = \bigl(\begin{smallmatrix}2
& 1\\ 0 & 3\end{smallmatrix}\bigr)$, $D_1 = 1$, $D_2 = 6$: $\Z^2/L
\cong \Z/6\Z$, de cardinalidade $\abs{\det B} = 6$.
\end{solution}

\begin{solution}{exo:b3:modules:5}
(a) Submódulo: feito na~\cref{def:b3:modules:torsion}. Se $a(x +
T(M)) = 0$ em $M/T(M)$ com $a \neq 0$,
então $ax \in T(M)$: $bax
= 0$ para algum $b \neq 0$, e $ba \neq 0$ (domínio), de modo que $x \in
T(M)$: a
classe é nula. $M/T(M)$ é sem \omterm{def:b3:modules:torsion}{torção}.

(b) $(X, Y) \subseteq K[X,Y]$ é sem \omterm{def:b3:modules:torsion}{torção} (submódulo do domínio $K[X,Y]$ agindo sobre si
mesmo). Suponha que fosse \omterm{def:b3:modules:free}{livre}; dois elementos $P, Q$ quaisquer
satisfazem $Q\cdot P - P \cdot Q = 0$, uma relação não trivial quando $P \ne Q$ são não nulos, de
modo que uma base tem um único elemento: $(X, Y) = (P)$ principal ---
contradizendo o~\cref{exo:b3:rings:6}(a). Sem \omterm{def:b3:modules:torsion}{torção} e \omterm{def:b3:modules:free}{finitamente gerado}
($X, Y$ geram) e, ainda assim, não \omterm{def:b3:modules:free}{livre}: sobre o não \omterm{def:b3:rings:pidufd}{DIP} $K[X,Y]$,
o teorema de estrutura falha.
\end{solution}

\begin{solution}{exo:b3:modules:6}
(a) Se $\Z = 2\Z \oplus C$, a projeção $\Z \to \Z/2\Z$ se restringe a um isomorfismo $C \cong \Z/2\Z$: $C$
seria um subgrupo de $\Z$ cujo elemento não nulo $x$ satisfaz
$2x \in C
\cap 2\Z = 0$. Mas $\Z$ é sem \omterm{def:b3:modules:torsion}{torção}: $C = 0$, forçando $\Z =
2\Z$ --- falso.

(b) $A^n/M$ é \omterm{def:b3:modules:free}{finitamente gerado} e sem \omterm{def:b3:modules:torsion}{torção}, logo \omterm{def:b3:modules:free}{livre}
(\cref{thm:b3:modules:structure}): $A^n/M \cong A^r$ com base $f_1, \dots, f_r$. Escolha pré-imagens $y_i \in A^n$ dos
$f_i$ e ponha $F = Ay_1 + \dots + Ay_r$. Todo $x \in A^n$ tem $\pi(x)
= \sum a_if_i$, de modo que $x - \sum a_iy_i \in M$: $A^n = M + F$.
Se $\sum a_iy_i \in M$, aplicar $\pi$ dá $\sum a_if_i = 0$, logo todos os $a_i = 0$ (base): $M \cap F = 0$.
Assim $A^n = M \oplus
F$.
\end{solution}

\begin{solution}{exo:b3:modules:7}
A redução do~\cref{exo:b3:modules:3} foi efetiva: com
\[
L = \begin{pmatrix} 1 & 0\\ -3 & 1\end{pmatrix},
\qquad
R = \begin{pmatrix} 1 & 2\\ 0 & -1 \end{pmatrix},
\qquad
LBR = \begin{pmatrix} 2 & 0\\ 0 & 4\end{pmatrix}
\]
(operação de linha $R_2 \leftarrow R_2 - 3R_1$, operações de coluna $C_2 \leftarrow C_2 - 2C_1$ e depois $C_2 \leftarrow -C_2$). Pondo
$y = R^{-1}x$ (uma bijeção de $(\Z/20\Z)^2$, pois $R$ é inversível sobre $\Z$), o
sistema $Bx \equiv b \pmod{20}$ é equivalente a
\[
2y_1 \equiv b_1, \qquad 4y_2 \equiv -3b_1 + b_2 \pmod{20}.
\]
A congruência $ky \equiv c \pmod{20}$ é \omterm{def:b3:groups:derived}{solúvel} se, e somente se, $\gcd(k, 20) \mid c$: há soluções se, e
somente se, $2 \mid b_1$ e $4
\mid b_2 - 3b_1$, isto é, \emph{$b_1$ par e $b_2 \equiv 3b_1
\pmod 4$}. Quando
\omterm{def:b3:groups:derived}{solúvel}, há $2 \times 4 = 8$ soluções \omterm{def:b3:modules:module}{módulo} $20$.
\end{solution}

\begin{solution}{exo:b3:modules:8}
(a) Uma $u$ nilpotente tem $\mu_u = X^k$; os \omterm{thm:b3:modules:structure}{divisores elementares} são $X^{k_1} \geq \dots$,
um bloco de Jordan $J_{k_i}(0)$ por parte de uma partição de $4$:
\begin{center}
\begin{tabular}{l l l}
partição & \omterm{thm:b3:modules:jordan}{forma de Jordan} & \omterm{thm:b3:modules:smith}{fatores invariantes}\\
$(4)$ & $J_4$ & $X^4$\\
$(3,1)$ & $J_3 \oplus J_1$ & $X,\ X^3$\\
$(2,2)$ & $J_2 \oplus J_2$ & $X^2,\ X^2$\\
$(2,1,1)$ & $J_2 \oplus J_1 \oplus J_1$ & $X,\ X,\ X^2$\\
$(1,1,1,1)$ & $0$ & $X, X, X, X$
\end{tabular}
\end{center}
(b) Tome $u = J_2\oplus J_2$ e $v = J_2 \oplus J_1 \oplus J_1$: ambas têm $\chi = X^4$, $\mu = X^2$, mas \omterm{thm:b3:modules:smith}{fatores invariantes}
diferentes --- não são semelhantes (\cref{thm:b3:modules:frobenius}); também é possível
comparar postos: $\operatorname{rk} u = 2 \neq 1 =
\operatorname{rk} v$. Para $n \leq 3$: $\chi$ e $\mu$ determinam,
para cada autovalor $\lambda$ (sobre um corpo de decomposição), o tamanho
total $m_\lambda \leq 3$ dos blocos de $\lambda$ e o maior bloco $r_\lambda$; e uma partição de
$m \leq 3$ fica determinada por sua maior parte ($m = 3, r = 2$ força $(2,1)$,
etc.). Logo os \omterm{thm:b3:modules:structure}{divisores elementares} coincidem, e o~\cref{cor:b3:modules:descent} faz a
semelhança descer ao corpo de base.
\end{solution}

\begin{solution}{exo:b3:modules:9}
(i)$\Rightarrow$(ii): se $V = K[u]x$, então $V \cong
K[X]/\operatorname{Ann}(x)$, e $\operatorname{Ann}(x) = (\mu_u)$ (um polinômio aniquila
$x$ se, e somente se, aniquila todo $V = K[u]x$, pois $P(u)Q(u)x = Q(u)P(u)x$). Logo $n = \dim V = \deg \mu_u$;
como $\mu_u \mid \chi_u$ e $\deg\chi_u = n$, a monicidade dá $\mu_u =
\chi_u$.

(ii)$\Rightarrow$(iii): $\deg\chi_u = \sum_i \deg P_i$ e $\mu_u = P_s$ (\cref{thm:b3:modules:frobenius}); a igualdade dos graus
força $s = 1$.

(iii)$\Rightarrow$(i): $V \cong K[X]/(P_1)$ é cíclico, gerado pela pré-imagem de $\bar 1$.

Uma \omterm{def:b3:modules:companion}{matriz companheira} é o próprio caso $V = K[X]/(P)$: $x = \bar
1$, isto é, $e_1$,
é cíclico. Para uma matriz diagonal $\operatorname{diag}(\lambda_1, \dots, \lambda_n)$: $\chi = \prod
(X - \lambda_i)$, $\mu = \prod_{\lambda \text{ distintos}} (X -
\lambda)$; eles coincidem
se, e somente se, os $\lambda_i$ são dois a dois distintos: uma matriz diagonal
tem vetor cíclico se, e somente se, suas entradas diagonais são dois a
dois distintas (e então $x = (1, \dots, 1)$ serve: Vandermonde).
\end{solution}

\begin{solution}{exo:b3:modules:10}
Smith: $M = Q\operatorname{diag}(d_1,\dots,d_n)P$; o~\cref{exo:b3:modules:4} dá $[\Z^n : M\Z^n] = \prod\abs{d_i} =
\abs{\det M}$. Se $\det M = \pm 1$: a fórmula da adjunta
$M^{-1} = (\det M)^{-1}\,{}^{t}\!\operatorname{com}(M)$ tem entradas inteiras, de modo que $M \in GL_n(\Z)$; reciprocamente, $MM^{-1} = I$ dá
$\det M \cdot \det M^{-1} = 1$ em $\Z$, de modo que $\det M = \pm1$.

Subgrupos de índice $m$ em $\Z^2$: um tal subgrupo $L$ tem
posto $2$ (índice finito) e uma única base em \emph{forma normal de
Hermite} $\bigl(\begin{smallmatrix} a & b\\ 0 &
d\end{smallmatrix}\bigr)$: $d$ é caracterizado por $L \cap (\{0\}
\times \Z) = \{0\} \times d\Z$, $a$ por $\pi_1(L) = a\Z$
(primeiras coordenadas), e $b$ é então único \omterm{def:b3:modules:module}{módulo} $d$;
normalize $a, d
> 0$ e $0 \leq b < d$. O índice é $ad = m$. Contagem: para cada
divisor $d \mid m$ ($a = m/d$), há $d$ escolhas de $b$: total $\sum_{d \mid m} d = \sigma_1(m)$.
\end{solution}

\begin{solution}{exo:b3:modules:11}
(a) Pelo teorema de estrutura, $G \cong \prod_i\Z/d_i\Z$, e $x^k = e$ se desacopla coordenada a
coordenada. Em $\Z/d\Z$: $kx \equiv 0
\pmod d$ tem exatamente $\gcd(k, d)$ soluções
($x$ tem de ser múltiplo de $d/\gcd(k,d)$, e há $\gcd(k, d)$ desses). Multiplique
sobre os fatores.

(b) Se $G = \Z/d_s\Z$ é cíclico ($s = 1$), a contagem é $\gcd(k, d_s) \leq k$. Se $s \geq 2$: tome
$k = d_1$; a contagem é $\prod_i\gcd(d_1, d_i) = d_1^{\,s} > d_1$ (cada $\gcd$ vale $d_1$ pela cadeia de
divisibilidade): a equação $x^{d_1} = e$ tem mais de $d_1$ soluções.

(c) Em um corpo, $X^k - 1$ tem no máximo $k$ raízes (\cref{ch:b3:rings}: um
polinômio não nulo de grau $k$ sobre um domínio), de modo que todo
subgrupo finito $G \leq K^\times$ satisfaz o critério de (b): $G$ é cíclico --- a
demonstração estrutural de uma linha da ciclicidade de $\mathbb F_q^\times$,
complementando a demonstração por contagem do~\cref{ch:b3:galois}.
\end{solution}

\begin{solution}{exo:b3:modules:12}
(a) Se $MN = I$ com $N$ inteira: $\det M\det N = 1$ com ambos inteiros, de modo
que $\det M = \pm1$. Reciprocamente, se $\det M =
\pm1$, a fórmula dos cofatores $M^{-1} =
\frac1{\det M}\,{}^t\!\operatorname{com}(M)$ tem
entradas inteiras.

(b) A multiplicação à esquerda por $E^{-k}$ subtrai $k$ vezes a
linha $2$ da linha $1$; por $F^{-k}$, $k$ vezes a linha
$1$ da linha $2$. Dada $M = \bigl(\begin{smallmatrix}a & *\\ c &
*\end{smallmatrix}\bigr) \in SL_2(\Z)$, a primeira coluna $(a, c)$ é um
vetor unimodular ($\gcd(a, c) = 1$: ele divide $\det M = 1$). Rode Euclides em $(a, c)$
com essas operações de linha: após finitos passos, a coluna se torna
$(\pm1, 0)$. A matriz é agora $\pm
\bigl(\begin{smallmatrix}1 & b\\ 0 & 1\end{smallmatrix}
\bigr) = \pm E^{b}$ (o determinante permaneceu $1$). Resta
escrever $-I$ nos geradores: $(EF^{-1}E)^2 =
\bigl(\begin{smallmatrix}0 & 1\\ -1 &
0\end{smallmatrix}\bigr)^2 = -I$ (verifique o quadrado da matriz de
rotação). Desenrolando, $M$ é uma palavra em $E^{\pm1},
F^{\pm1}$.

(c) O algoritmo de Smith (\cref{met:b3:modules:smith}) usa exatamente essas operações
de linha e de coluna (mais trocas e mudanças de sinal, elas próprias
produtos de operações elementares a menos do sinal do determinante):
sobre $\Z$, toda matriz é $U\,D\,V$ com $U, V$ produtos de matrizes
elementares e $D$ a forma de Smith --- (b) é a instância $2\times2$, de
determinante $1$, do fato geral de que as matrizes do tipo $E$
geram $SL_n(\Z)$.
\end{solution}

\begin{solution}{pb:b3:modules:1}
\textbf{1.} $\varphi$ é aditiva e $K[X]$-linear: $\varphi(X
\cdot (Q_i)_i) = \sum_i (XQ_i)(M)e_i = M\sum_i Q_i(M)e_i = X
\cdot \varphi\bigl((Q_i)_i\bigr)$, sendo a estrutura
de \omterm{def:b3:modules:module}{módulo} de $V_M$ dada por $X \cdot v = Mv$. Ela é sobrejetora: os vetores
constantes dão todo $K^n$. A coluna $j$ de $XI - M$ é $Xe_j - \sum_i
m_{ij}e_i$,
cuja imagem é $Me_j - \sum_i m_{ij}e_i = 0$.

\textbf{2.} \omterm{def:b3:modules:module}{Módulo} as colunas, $Xe_j \equiv \sum_i m_{ij}e_i$: todo vetor de polinômios se reduz,
por indução sobre o grau máximo, a um vetor \emph{constante} $c \in K^n$. Se o
vetor original está em $\ker\varphi$, então $\varphi(c) = c = 0$ (nas constantes, $\varphi$ é a
identificação $K^n = V_M$), de modo que o vetor está no espaço gerado pelas
colunas: $\ker\varphi = (XI_n -
M)K[X]^n$. Logo $V_M \cong K[X]^n/(XI - M)K[X]^n$, e Smith sobre $K[X]$ (todos os \omterm{thm:b3:modules:smith}{fatores
invariantes} não nulos, com produto $\det(XI - M) = \chi_M$) dá $V_M \cong \bigoplus_i
K[X]/(f_i)$: os $f_i$ não
constantes são os \omterm{thm:b3:modules:frobenius}{invariantes de semelhança}, calculáveis como quocientes
de mdc de menores (\cref{thm:b3:modules:smith}).

\textbf{3.} $\lambda I_n$: $XI - \lambda I$ já está em forma de Smith: invariantes $(X - \lambda, \dots, X - \lambda)$,
$n$ deles. Entradas diagonais distintas: $V \cong \bigoplus_i K[X]/(X -
\lambda_i)$ com \omterm{def:b3:modules:module}{módulos} dois a
dois \omterm{thm:b3:rings:crt}{comaximais}, de modo que o teorema chinês dos restos comprime tudo no
único cíclico $K[X]/\bigl(\prod_i(X - \lambda_i)\bigr)$: um invariante, $\chi$. $J_n(0)$: $\mu = X^n = \chi$ força
um único invariante $X^n$. $\operatorname{diag}(J_2(0), J_1(0))$: \omterm{thm:b3:modules:structure}{divisores elementares} $X^2, X$:
invariantes $P_1 = X \mid P_2 =
X^2$.

\textbf{4.} $P \mapsto P(u)$ leva $K[X]$ sobre $K[u]$, com núcleo $(\mu_u)$ pela
definição de polinômio minimal: $K[u]
\cong K[X]/(\mu_u)$, de dimensão $\deg\mu_u = \deg P_s = n_s$.

\textbf{5.} A $K[X]$-linearidade força $f(\bar Q) = f(Q\cdot\bar 1)
= Q\,f(\bar 1)$. A classe $\bar 1$
satisfaz $P \bar 1 = 0$, de modo que $P f(\bar 1) = 0$ é necessário. Reciprocamente, se $P\bar c = 0$,
então $f(\bar Q) = Q\bar c$ está bem definida ($Q \equiv Q' \bmod P
\Rightarrow (Q - Q')\bar c \in$ múltiplos de $P\bar c = 0$) e é
$K[X]$-linear.

\textbf{6.} Sejam $g = \gcd(P, Q)$, $P = gP'$, $Q = gQ'$ com $\gcd(P', Q') = 1$. Em $K[X]/(Q)$: $P\bar c = 0 \iff Q \mid Pc
\iff Q' \mid P'c \iff Q' \mid c$
(Euclides, $\gcd(P', Q') = 1$). Logo os $\bar c$ admissíveis formam o submódulo
gerado por $\bar{Q'}$, cujo anulador é $\{R : Q \mid RQ'\} = (g)$: esse submódulo é $\cong K[X]/(g)$. Com a
questão 5, $\operatorname{Hom}(K[X]/(P), K[X]/(Q)) \cong K[X]/(\gcd(P,Q))$, de dimensão $\deg\gcd(P,Q)$.

\textbf{7.} $v$ comuta com $u$ se, e somente se, $v$ comuta
com todo $P(u)$, se, e somente se, $v$ é $K[X]$-linear:
$\mathcal C(u) =
\operatorname{End}_{K[X]}(V)$. Escrevendo os morfismos de $V =
\bigoplus_j V_j$ como matrizes $(f_{ij})$, $f_{ij} \in
\operatorname{Hom}(V_j, V_i)$
(componha com injeções e projeções), a questão 6 dá
\[
\dim \mathcal C(u) = \sum_{i,j} \deg\gcd(P_i, P_j)
= \sum_{i,j} n_{\min(i,j)}
= \sum_{k=1}^{s} \bigl(2(s - k) + 1\bigr)\,n_k ,
\]
usando a cadeia de divisibilidade ($\gcd(P_i, P_j) =
P_{\min(i,j)}$) e, na última passagem, que
$\min(i,j) = k$ ocorre para exatamente $2(s - k) + 1$ pares $(i, j)$.

\textbf{8.} Como $2(s-k)+1 \geq 1$, $\dim\mathcal C(u) \geq
\sum_k n_k = n$, com igualdade se, e somente se, $s = 1$,
isto é, se, e somente se, $u$ é cíclico (\cref{exo:b3:modules:9}). Para $u = \lambda\,\mathrm{id}$:
$s = n$, todos $n_k = 1$: $\dim = \sum_{k=1}^n (2(n-k)+1) = n^2$ --- correto, pois $\mathcal C(\lambda\,\mathrm{id}) = \mathcal
L(V)$.

\textbf{9.} Pela fórmula: invariantes $(X, X^2)$, de modo que $s = 2$,
$n_1 =
1$, $n_2 = 2$: $\dim = 3\cdot1 + 1\cdot2 = 5$. Diretamente: na base $(e_1, e_2, e_3)$ com $u e_2 = e_1$, $ue_1 = ue_3 = 0$, escrever
$Au = uA$ para $A = (a_{ij})$ dá as condições $a_{21} = a_{23} = a_{31} = 0$ e $a_{11} = a_{22}$: cinco parâmetros \omterm{def:b3:modules:free}{livres}
$a_{11}{=}a_{22}, a_{12}, a_{13}, a_{32}, a_{33}$.

\textbf{10.} Um polinômio $P(u)$ comuta com tudo o que comuta com
$u$ (ele é uma soma de potências de $u$): $K[u]
\subseteq \mathcal C(\mathcal C(u))$. E $u \in \mathcal C(u)$, de modo
que todo $w \in \mathcal C(\mathcal C(u))$ comuta com $u$.

\textbf{11.} Sejam $V = K[u]x$ e $v \in \mathcal C(u)$. Escreva $v(x) = P(u)x$ (ciclicidade). Para $y = Q(u)x$
arbitrário: $v(y) =
vQ(u)x = Q(u)v(x) = Q(u)P(u)x = P(u)y$. Logo $v = P(u)$: $\mathcal
C(u) = K[u]$. Então $\mathcal C(\mathcal C(u)) = \mathcal
C(K[u]) = \mathcal C(u) = K[u]$ (comutar com todo
$K[u]$ é o mesmo que comutar com $u$). O teorema vale no caso
cíclico.

\textbf{12.} $\pi_i$ é $K[X]$-linear (a decomposição é uma soma
direta de submódulos), de modo que $\pi_i \in \mathcal C(u)$, e $w$ comuta com ela:
$w(V_i) = w\pi_i(V) = \pi_i w(V) \subseteq
V_i$. A restrição $w_i = w\restriction_{V_i}$ comuta com o cíclico $u_i = u\restriction_{V_i}$ (questão 10), e $w_i
\in \mathcal C(u_i) = K[u_i]$
(questão 11): $w_i = Q_i(u_i) =
Q_i(u)\restriction_{V_i}$ para algum $Q_i \in K[X]$.

\textbf{13.} Boa definição de $\eta_{ij}(R(u)x_j) = R(u)x_i$: se $R(u)x_j = 0$, então $P_j \mid R$, e $P_i \mid P_j$ dá
$P_i \mid R$, de modo que $R(u)x_i = 0$ ($\operatorname{Ann}(x_i) =
(P_i)$). $\eta_{ij}$ é então $K[X]$-linear por
construção, e $\tilde\eta_{ij} = \eta_{ij}\pi_j$ é uma composição de $K[X]$-morfismos $V \to V$: $\tilde\eta_{ij} \in \mathcal
C(u)$.

\textbf{14.} Avalie $w\tilde\eta_{ij} = \tilde\eta_{ij}w$ em $x_j$: o membro esquerdo é $w(x_i) = Q_i(u)x_i$; o
direito é $\eta_{ij}\bigl(Q_j(u)x_j\bigr) = Q_j(u)x_i$. Logo $(Q_i -
Q_j)(u)\,x_i = 0$: $P_i \mid Q_i - Q_j$ para todo $i \leq j$. Em particular, com
$Q = Q_s$: $Q \equiv Q_i \pmod{P_i}$, de modo que $Q(u)$ e $Q_i(u)$ coincidem em $V_i$
(que $P_i(u)$ aniquila). Portanto $w
= Q(u)$ em cada $V_i$ e, assim,
em $V$: $\mathcal C(\mathcal C(u))
\subseteq K[u]$ e, com a questão 10, $\mathcal C(\mathcal
C(u)) = K[u]$.

\textbf{15.} A primeira afirmação são as questões 10--14 (ou, para
$u$ cíclico, apenas a questão 11). Para $u = \mathrm{id}$, $n \geq
2$: $\mathcal C(u) = \mathcal L(V)$ tem
dimensão $n^2$, ao passo que $\dim K[u] = \deg\mu_u = 1$. Logo $\mathcal C(u) = K[u]$ falha
espetacularmente; e, no entanto, o teorema do bicomutante vale ($\mathcal C(\mathcal L(V)) = K\,\mathrm{id} = K[u]$: o
\omterm{ex:b3:groups:actions}{centro} da álgebra de matrizes são os escalares). A ciclicidade é o que
torna o \emph{primeiro} comutante já polinomial; o comutante
\emph{duplo} é polinomial sempre.

\textbf{16.} Se $P(XI - M)Q = S$ é uma redução de Smith (com $P, Q$ inversíveis
sobre $K[X]$), transpor dá ${}^tQ\,(XI -
{}^tM)\,{}^tP = {}^tS = S$: mesma forma de Smith, de modo que
$XI - M$ e $XI - {}^tM$ têm os mesmos \omterm{thm:b3:modules:smith}{fatores invariantes}, isto é, $M$ e
${}^tM$ têm os mesmos \omterm{thm:b3:modules:frobenius}{invariantes de semelhança}
(\cref{cor:b3:modules:descent}): elas são semelhantes.

\textbf{17.} Os \omterm{thm:b3:modules:frobenius}{invariantes de semelhança} de $M$ sobre $L$ são os
\omterm{thm:b3:modules:smith}{fatores invariantes} de $XI - M$ em $L[X]$. Uma redução de Smith
de $XI - M$ sobre $K[X]$ --- com $P, Q$ inversíveis sobre
$K[X]$, diagonal e com a cadeia de divisibilidade --- é
\emph{também} uma redução de Smith válida sobre $L[X]$ ($P, Q$
permanecem inversíveis: seus determinantes são constantes não nulas), e
os \omterm{thm:b3:modules:smith}{fatores invariantes} mônicos são únicos: os \omterm{thm:b3:modules:smith}{fatores invariantes}
calculados sobre $K$ e sobre $L$ coincidem. Logo $M \sim_L N$ se, e
somente se, elas têm os mesmos \omterm{thm:b3:modules:smith}{fatores invariantes}, se, e somente se,
$M \sim_K N$. Em particular, matrizes reais $\C$-conjugadas são
$\R$-conjugadas --- um enunciado muitas vezes demonstrado por via
analítica (especialize um $P + \iu
Q$ inversível), aqui obtido
estruturalmente.

\textbf{18.} Decomponha $u = \bigoplus J_{m_t}(0)$ em blocos de Jordan nilpotentes. Em um
bloco de tamanho $m$, $\operatorname{rk}J_m^k = \max(m - k, 0)$, de modo que $\operatorname{rk}J_m^{k-1} - \operatorname{rk}J_m^k = 1$ se $m
\geq k$, e $0$
caso contrário. Somando sobre os blocos: $\operatorname{rk}u^{k-1} - \operatorname{rk}u^k =
\#\{t : m_t \geq k\}$. A sequência dos postos
determina, portanto, o multiconjunto $(m_t)$ --- uma partição de
$n$ --- e, reciprocamente, toda partição é realizada: as classes de
nilpotentes $\leftrightarrow$ as partições de $n$, sobre qualquer corpo. $M_5$:
$p(5) = 7$ classes ($5$; $4{+}1$; $3{+}2$; $3{+}1{+}1$; $2{+}2{+}1$; $2{+}1{+}1{+}1$;
$1^5$).

\textbf{19.} (a) Sobre $\Q$ (ou qualquer corpo de característica
$0$), $D^n = 0$, $D^{n-1}(X^{n-1}) = (n-1)!\, \neq 0$: $D$ é nilpotente de índice $n$ em um
espaço de dimensão $n$, logo cíclico com único invariante $X^n$
($x = X^{n-1}$ gera: suas derivadas iteradas geram o espaço). (b) Sobre $\mathbb F_p$ com
$p
< n$: $D^p = 0$, pois a $p$-ésima derivada de todo monômio
$X^m$ carrega o fator $m(m-1)\cdots(m-p+1)$, um produto de $p$ inteiros
consecutivos, logo $\equiv 0 \bmod p$. Escreva $n = ap + r$, $0 \leq r < p$. Então $\ker D^k$ é gerado pelos
monômios $X^m$ com $D^kX^m = 0$; contando os expoentes $m < n$ por seu
resto \omterm{def:b3:modules:module}{módulo} $p$: $\dim\ker D^k = ak + \min(r,
k)$ para $0 \leq k \leq p$, de modo que $\operatorname{rk}D^{k-1} -
\operatorname{rk}D^k = \dim\ker D^k - \dim\ker D^{k-1} = a +
\mathbf 1_{k \leq r}$. Pela questão
18, a partição tem $a$ blocos de tamanho exatamente $p$ e (se
$r > 0$) um bloco de tamanho $r$: \omterm{thm:b3:modules:frobenius}{invariantes de semelhança}
$X^r \mid X^p \mid \dots \mid X^p$. A característica muda a forma canônica do operador mais familiar
da matemática.

\textbf{20.} $M \in GL_2$ tem $s \in \{1, 2\}$ \omterm{thm:b3:modules:smith}{fatores invariantes}. Se $s = 2$: $P_1 = P_2 = X - a$
($a \neq 0$: inversibilidade), isto é, $M = aI$, central. Se $s = 1$:
$M$ é cíclica, com polinômio característico $=$ minimal
$\chi$ de grau $2$, e as classes correspondem aos $\chi$
possíveis com $\chi(0) \neq 0$: $\chi$ decomposto com raízes distintas $a \neq
b$
(companheira $\sim$ diagonal); $\chi = (X - a)^2$ (companheira, não semissimples);
$\chi$ \omterm{def:b3:rings:divisibility}{irredutível}. Exatamente um tipo em cada caso --- os \omterm{thm:b3:modules:smith}{fatores
invariantes} são um invariante completo.

\textbf{21.} Central: $q - 1$ escolhas de $a$. Autovalores
decompostos distintos: pares não ordenados $\{a, b\} \subseteq
\mathbb F_q^\times$, $a \neq b$: $\binom{q-1}2$ classes.
Minimal $(X-a)^2$: $q - 1$ classes. Quadráticas irredutíveis com termo
constante não nulo: todas as quadráticas irredutíveis servem (suas
raízes são não nulas), e há $\frac{q^2 - q}2$ quadráticas mônicas irredutíveis (as
$q^2$ quadráticas mônicas menos as $\binom q2 + q = \frac{q^2+q}2$ decompostas). Total:
\[
(q - 1) + \frac{(q-1)(q-2)}2 + (q - 1) + \frac{q^2 - q}2
= q^2 - 1 .
\]

\textbf{22.} Se $M$ não é cíclica, $s = 2$ e $M$ é escalar (a
dicotomia da questão 20 vale em $M_2$, inversível ou não: dois
\omterm{thm:b3:modules:smith}{fatores invariantes} de grau $1$ com $P_1 \mid P_2$ e $P_1 =
P_2$ forçam $M = aI$).
Os escalares são em número de $q$ entre as $q^4$ matrizes:
probabilidade de ser cíclica $1 - q^{-3}$. Em geral, o lugar não cíclico de
$M_n$ é onde os menores $(n-1)\times(n-1)$ de $XI - M$ compartilham um fator
--- uma condição algébrica própria ---, de modo que sua proporção é
pequena, $O(1/q)$, para $q$ grande: as matrizes com $\chi = \mu$ são a
regra, e o argumento de colagem da Parte III é o preço pago pelas
exceções.

\textbf{23.} Um elemento do \omterm{ex:b3:groups:actions}{centro} de $\mathcal C(u)$ está em $\mathcal C(u)$ e comuta com
todo elemento de $\mathcal C(u)$, isto é, está em $\mathcal C(\mathcal C(u)) =
K[u]$ (questão 14).
Reciprocamente, $K[u] \subseteq \mathcal C(u)$, e todo $P(u)$ comuta com todo $v \in \mathcal C(u)$ (um tal
$v$ comuta com $u$, logo com cada potência de $u$):
$K[u]$ é central em $\mathcal C(u)$. Logo $Z(\mathcal C(u))
= K[u]$. Consequentemente, $\mathcal C(u)$ é
comutativa se, e somente se, $\mathcal C(u) = Z(\mathcal C(u)) = K[u]$; nesse caso $\dim\mathcal C(u) = \dim K[u] = n_s \leq n$, ao passo que a
questão 8 dá $\dim\mathcal C(u) \geq n$: logo $n_s = n$ e $s = 1$, isto é, $u$ é cíclico.
Reciprocamente, para $u$ cíclico a questão 11 dá $\mathcal C(u) = K[u]$, comutativa.
Estruturalmente: uma álgebra de matrizes igual a seu próprio \omterm{ex:b3:groups:actions}{centro} é
exatamente uma álgebra de polinômios $K[u]$ de um $u$ cíclico.

\textbf{24.} Cada coeficiente $2s - 2i + 1$ na fórmula de Frobenius é ímpar, de
modo que
\[
\dim\mathcal C(u) = \sum_{i=1}^s(2s - 2i + 1)\,n_i
\equiv \sum_{i=1}^s n_i = n \pmod 2 .
\]
Para $n = 4$, as sequências de graus $n_1 \leq \dots
\leq n_s$ possíveis dos \omterm{thm:b3:modules:smith}{fatores
invariantes}, de soma $4$, são $(4)$, $(1, 3)$, $(2, 2)$,
$(1, 1, 2)$, $(1, 1, 1, 1)$; cada uma é realizada, por exemplo pelo nilpotente com $P_i = X^{n_i}$
(a cadeia de divisibilidade vale automaticamente). A fórmula dá,
respectivamente,
\[
1\cdot4 = 4, \quad 3 + 3 = 6, \quad 6 + 2 = 8, \quad
5 + 3 + 2 = 10, \quad 7 + 5 + 3 + 1 = 16 .
\]
Logo o conjunto de valores é $\{4, 6, 8, 10, 16\}$: números pares da paridade certa, mas
$12$ e $14$ jamais ocorrem --- entre as sequências quase cíclicas
e o $n^2$ do escalar há uma lacuna.

\textbf{25.} $\abs{GL_2(\mathbb F_3)} = (q^2 - 1)(q^2 - q) =
8 \cdot 6 = 48$. Tipo central: $I$ e $2I$, duas classes de
tamanho $1$. Nos três outros tipos, $M$ é cíclica (questão 20), de
modo que seu centralizador em $GL_2$ é o grupo das unidades de $\mathcal C(M) = K[M]$
(questão 11), e o tamanho da classe é $=48/\abs{K[M]^\times}$ por órbita--estabilizador.
Autovalores decompostos distintos: apenas o par $\{1, 2\}$, uma classe; $K[M]
\cong \mathbb F_3 \times \mathbb F_3$
(teorema chinês dos restos em $\chi = (X-1)
(X-2)$), unidades $2 \cdot 2 = 4$, tamanho $48/4 = 12$.
Minimal $(X - a)^2$, $a \in \{1, 2\}$: duas classes; $K[M] \cong
\mathbb F_3[X]/((X-a)^2)$, unidades $q^2 - q = 6$ (o termo
constante da unidade é $\neq 0$ após centrar), tamanho $48/6 = 8$.
$\chi$ \omterm{def:b3:rings:divisibility}{irredutível}: as quadráticas mônicas irredutíveis sobre $\mathbb F_3$
são em número de $(9 - 3)/2 = 3$, a saber
\[
X^2 + 1, \qquad X^2 + X + 2, \qquad X^2 + 2X + 2,
\]
(sem raízes em $\mathbb F_3$: verifique $0, 1, 2$); três classes, $K[M] \cong \mathbb F_9$, unidades
$q^2 - 1 = 8$, tamanho $48/8 =
6$. \omterm{cor:b3:groups:classeq}{Equação das classes}: $2\cdot1 + 1\cdot12 + 2\cdot8 + 3\cdot6 =
2 + 12 + 16 + 18 = 48$; e $2 + 1 + 2 + 3 = 8 = q^2 - 1$ classes,
coincidindo com a questão 21.
\end{solution}
